\documentclass[11pt]{article}
\usepackage[T1]{fontenc}
\usepackage{lmodern,amsmath,amssymb,amsthm,mathtools,microtype}
\usepackage[margin=1in]{geometry}
\usepackage[hidelinks]{hyperref}
\newtheorem{theorem}{Theorem}[section]
\newtheorem{lemma}[theorem]{Lemma}
\newtheorem{proposition}[theorem]{Proposition}
\newtheorem{corollary}[theorem]{Corollary}
\newcommand{\R}{\mathbb R}\newcommand{\C}{\mathbb C}
\newcommand{\E}{\mathbb E}\newcommand{\norm}[1]{\left\lVert#1\right\rVert}
\newcommand{\code}[1]{{\small\nolinkurl{#1}}}
\title{Classical strip contraction and an attained threshold\\for the actual quadratic Xi flow}
\author{Project Stardust\\\small A classical formalization companion}
\date{Reviewed editorial draft, October 30, 2026}
\begin{document}\maketitle
\begin{abstract}
We reconstruct the classical coarse zero-strip bound for the quadratic
heat deformation of the actual arithmetic Xi kernel. For the normalization
$F_t(z)=\int_0^\infty K(u)e^{tu^2}\cos(zu)\,du$, every zero at
$t\ge0$ satisfies $(\operatorname{Im}z)^2\le\max(1-2t,0)$.
In particular all zeros are real for $t\ge1/2$.
The proof constructs real even divisor polynomials retaining the initial
zeros and one arbitrarily small quadratic growth budget shared by all
approximants. This budget justifies locally uniform imaginary-Gaussian
heat convergence. A conjugate-root product identity contracts zero strips
under symmetric imaginary shifts, and a finite-degree Euler limit gives
the actual polynomial heat operator. A vertical modulus argument transfers
the strip even across repeated roots and collisions. We include the exact
Xi frequency normalization, positive origin value and original Gaussian
integrability before passing to limits. These are classical de Bruijn
mechanisms, already stated in broader form by Newman; no new strip bound,
exact threshold, zero-time reality or Riemann-hypothesis result is claimed.
Exact declarations and hypotheses are mapped to a frozen Lean library.
The same actual flow has a finite attained all-real threshold; a complete
four-point smoothing proof supplies the needed lower obstruction.
\end{abstract}

\section{The arithmetic flow and the conclusion}
For $v>0$ write $\Theta_k(v)=\sum_{n\ge1}n^ke^{-\pi n^2v}$ and, for
$u\ge0$, set
\begin{equation}\label{eq:kernel}
 K(u)=2\pi^2e^{9u}\Theta_4(e^{4u})-3\pi e^{5u}\Theta_2(e^{4u}).
\end{equation}
The quadratic flow is the actual absolutely convergent integral
\begin{equation}\label{eq:flow}
 F_t(z)=\int_0^\infty K(u)e^{tu^2}\cos(zu)\,du,
       \qquad t\in\R,\ z\in\C.
\end{equation}
Time has the \emph{growing} sign $e^{tu^2}$, so its equation is
$\partial_tF_t=-\partial_z^2F_t$. It is distinct from the positive-time
Fourier multiplier $e^{-tu^2}$ of ordinary real heat evolution.
Let
\[
 \xi(s)=\tfrac12s(s-1)\pi^{-s/2}\Gamma(s/2)\zeta(s)
\]
with its entire continuation at the apparent exceptional points. Our
frequency and scalar normalizations are
\begin{equation}\label{eq:xi}
 F_0(z)=\frac18\xi\left(\frac12+\frac{iz}{2}\right).
\end{equation}
Thus the elementary critical strip $0\le\operatorname{Re}s\le1$
corresponds to $|\operatorname{Im}z|\le1$, rather than width $1/2$.

\begin{theorem}\label{thm:main}
For every real $t\ge0$ and $z\in\C$,
\[
 F_t(z)=0\quad\Longrightarrow\quad
       (\operatorname{Im}z)^2\le\max(1-2t,0).
\]
In particular $F_t$ has only real zeros for every $t\ge1/2$.
Each $F_t$ is an even real entire function and $F_t(0)>0$.
\end{theorem}
The all-real conclusion includes multiplicities; it does not assert
simplicity. The upper bound $1/2$ is coarse and classical. No assertion
about the exact infimum of all-real times, or about zeros at time zero,
is made. The fixed mathematical baseline is
\begin{center}\small\texttt{bab688bb6140cd6316c80680095a914ea0d346db}.\end{center}

Newman's Theorem 1, printed page 246~\cite{Newman}, already states the
broader classical strip contraction for even rapidly decaying Fourier
kernels, crediting de Bruijn~\cite{deBruijn}. Our purpose is to expose how
the actual arithmetic and approximation hypotheses in the verified
library close that argument. This is a formalization exposition, not a
new analytic theorem.

\section{Integrability, positivity and the exact Xi normalization}
The theta series and their derivatives converge uniformly on compact
subsets of $v>0$. For example any polynomial factor $n^k$ is dominated
by a constant times $e^{\pi n^2a/2}$ on $v\ge a>0$, leaving a summable
Gaussian tail. Termwise differentiation consequently gives
$\Theta_k'=-\pi\Theta_{k+2}$ there.
Expanding \eqref{eq:kernel} yields
\[
 K(u)=\sum_{n\ge1}\pi n^2e^{5u}
       (2\pi n^2e^{4u}-3)e^{-\pi n^2e^{4u}}>0\qquad(u\ge0),
\]
since $2\pi>3$. Moreover, for $v\ge1$,
\[
 \Theta_k(v)\le\Theta_k(1/2)e^{-\pi v/2}.
\]
Indeed $n^2v\ge n^2/2+v/2$ for $n\ge1$. Thus for a finite constant $C$,
\begin{equation}\label{eq:tail}
 |K(u)|\le C e^{9u-\pi e^{4u}/2}\quad(u\ge0).
\end{equation}
Every fixed nonnegative integer power of $u$, multiplied by
$e^{Tu^a+Bu}$ with a fixed nonnegative integer $a$ and real $T,B$,
is integrable against this bound. This follows because $e^{4u}$
dominates every fixed power of $u$; beyond a finite cutoff the full
integrand is bounded by a constant times $e^{-u}$. On the remaining
compact interval continuity supplies integrability.

For $z$ in any fixed complex disk, $|\cos(zu)|\le e^{|z|u}$, and each
spatial derivative adds only a power of $u$. The same bound with a larger
disk is a common integrable majorant. Differentiation under the integral
therefore proves entireness, all spatial derivatives and the stated time
equation. Evenness is literal; conjugation gives
$F_t(\overline z)=\overline{F_t(z)}$. At zero the integrand is strictly
positive, so $F_t(0)>0$. In particular none of the functions to which the
zero-limit argument will be applied is identically zero.

For completeness the normalization can be recovered without treating
\eqref{eq:xi} as a new assumption about the kernel. Put
\[
 \psi(u)=e^u\Theta_0(e^{4u}).
\]
Direct differentiation gives $K=(\psi''-\psi)/8$. Jacobi's classical
theta identity implies
\[
 \psi(u)-\psi(-u)=\frac{e^{-u}-e^u}{2},\qquad \psi'(0)=-\frac12.
\]
The derivative follows by differentiating the first identity at zero.
The functions $\psi,\psi'$ decay at positive infinity with the same
superexponential mechanism as \eqref{eq:tail}. For real $x$ the derivative
of
$\psi'(u)\cos(xu)+x\psi(u)\sin(xu)$ is
$(\psi''(u)+x^2\psi(u))\cos(xu)$. Integrating and evaluating both boundary
terms gives
\begin{equation}\label{eq:primitive}
 F_0(x)=\frac1{16}-\frac{1+x^2}{8}
                    \int_0^\infty\psi(u)\cos(xu)\,du.
\end{equation}
The standard theta Mellin representation of the pole-removed completed
zeta, split at one and transformed by the same Jacobi identity, is
\[
 \Lambda_0(s)=\int_1^\infty\Theta_0(v)
                 \bigl(v^{s/2}+v^{(1-s)/2}\bigr)\,\frac{dv}{v}.
\]
Here $\Lambda_0$ is the entire part after removing the two simple pole
terms, so $\xi(s)=(s(s-1)\Lambda_0(s)+1)/2$. The integral converges for
all $s$ by exponential theta decay, and is the usual analytic continuation
of its convergent Mellin representation. At $s=(1+ix)/2$, substitution
$v=e^{4u}$ gives
$\Lambda_0(s)=8\int_0^\infty\psi(u)\cos(xu)\,du$.
Since $s(s-1)=-(1+x^2)/4$, this identity and \eqref{eq:primitive} prove
\eqref{eq:xi} on the real axis. Both sides are entire, so the identity
theorem proves it on the entire complex plane. The foundational inputs
here are the classical Jacobi theta identity and Mellin representation;
the displayed derivatives, substitutions and boundary evaluations are
the actual normalization calculation in the source library.

\subsection*{The initial strip and a subquadratic bound}
The Euler product gives $\zeta(s)\ne0$ for $\operatorname{Re}s>1$.
Equation~\eqref{eq:xi}, the nonvanishing gamma factor and the functional
equation therefore exclude zeros with $\operatorname{Im}z<-1$ or
$\operatorname{Im}z>1$. Apparent exceptional points in the completed
formula do not introduce zeros: on the imaginary axis the actual integral
is $\int K(u)\cosh(yu)du>0$.
Thus
\begin{equation}\label{eq:initialstrip}
 F_0(z)=0\quad\Longrightarrow\quad |\operatorname{Im}z|\le1.
\end{equation}
No improved arithmetic zero-free half-plane is imported.

For fixed $t$, there is $C_t>0$ such that
\begin{equation}\label{eq:growth}
 |F_t(z)|\le C_t e^{|z|^{3/2}}\qquad(z\in\C).
\end{equation}
Here is an elementary bound matching the source exponent. If $R\ge0$
and $u\ge0$, then $Ru\le R^{3/2}+u^3$: for $u\le\sqrt R$ use
$Ru\le R^{3/2}$; otherwise use $Ru\le u^3$. Therefore
\[
 |F_t(z)|\le e^{|z|^{3/2}}\int_0^\infty
                      |K(u)|e^{tu^2+u^3}\,du,
\]
and the integral is finite by \eqref{eq:tail}. This is a genuine
subquadratic entire growth bound, not only a real-axis estimate.

\section{Real divisor polynomials with one common Gaussian budget}
We use the classical Hadamard factorization theorem under its explicit
finite-order hypotheses. An even entire $f$ with $f(0)\ne0$ and
$|f(z)|\le C e^{|z|^{3/2}}$ has the paired factorization
\begin{equation}\label{eq:product}
 f(z)=f(0)\prod_{\rho\in\mathcal Z_+}
                (1-z^2/\rho^2)^{m_\rho}.
\end{equation}
Here $\mathcal Z_+$ contains one representative of each nonzero opposite
zero pair, and $m_\rho$ is the actual multiplicity. Products are ordered
by expanding disks. Hadamard's genus-one elementary factors for opposite
zeros multiply to $1-z^2/\rho^2$; the residual exponential has degree at
most one, and evenness makes its linear coefficient zero. This is the
classical entire-function input, not an asserted Lean theorem in the
shorter notation above. McMullen's primary notes state the required
Hadamard theorem and opposite-pair simplification~\cite{McMullen}.

The inverse squares are absolutely summable. One way to see the exact
hypothesis is Jensen's formula: $f(0)\ne0$ and the growth bound give
$n_f(R)\le C'(1+R^{3/2})$ by applying Jensen at $2R$. On dyadic annuli,
\[
 \sum_{2^j\le|\rho|<2^{j+1}}\frac{m_\rho}{|\rho|^2}
 \le\frac{n_f(2^{j+1})}{2^{2j}}=O(2^{-j/2}),
\]
so summing the tail proves the assertion. There are finitely many zeros
in bounded disks and none in a sufficiently small disk around zero.
These facts justify uniform convergence of the paired product on compact
sets and all finite products used below.

Apply this to $f=F_0$. Choose radii $R_n\to\infty$ avoiding zeros, and let
\begin{equation}\label{eq:polys}
 P_n(z)=F_0(0)\prod_{\substack{\rho\in\mathcal Z_+\\|\rho|<R_n}}
                          (1-z^2/\rho^2)^{m_\rho}.
\end{equation}
Complete disks preserve all conjugate and opposite zeros, so $P_n$ is a
real even polynomial. It is nonzero because $P_n(0)=F_0(0)>0$.
Its zeros, with their multiplicities, are exactly the zeros of $F_0$ in
that disk, and all satisfy \eqref{eq:initialstrip}. Also $P_n\to F_0$
locally uniformly. The source makes a definite expanding-radius choice
$2(n+1)<R_n<2(n+1)+1$; no monotonicity of that particular choice is
needed. Radii can always be chosen away from the discrete zero set.

\begin{lemma}\label{lem:budget}
For every $\epsilon>0$ there is $C_\epsilon>0$, independent of $n$, with
\[
 |P_n(z)|\le C_\epsilon e^{\epsilon|z|^2}\quad(n\ge0,\ z\in\C).
\]
\end{lemma}
\begin{proof}
Choose a finite initial set of paired zeros so the sum of
$m_\rho/|\rho|^2$ outside it is at most $\epsilon/2$. For any finite
subset of that tail,
\[
 \prod|1-z^2/\rho^2|^{m_\rho}
 \le\exp\left(|z|^2\sum\frac{m_\rho}{|\rho|^2}\right)
 \le e^{\epsilon|z|^2/2}.
\]
For each of the finitely many initial factors, its polynomial growth is
bounded by a constant times an arbitrarily small Gaussian. More
explicitly, $1+a r^2\le(1+a/\delta)e^{\delta r^2}$ for $a\ge0$,
$\delta>0$. Choose $\delta$ so the sum of the allocated multiplicity
budgets is at most $\epsilon/2$. This bounds every subset of the finite
initial product by one constant times $e^{\epsilon|z|^2/2}$, enlarging
the constant to at least one for omitted factors. Combine head and tail
and include $|F_0(0)|$. The same bound handles the finitely many early
approximants and establishes a single constant for all $n$.
\end{proof}
The source proves this common budget by finite divisor jets and an
absolute annular tail estimate in the initial strip. The Jensen/Hadamard
presentation above is a labelled classical human reorganization of the
same actual divisor-polynomial conclusion. It does not assume that
$F_0$ has real or simple zeros.

\section{The Gaussian heat limit is in its actual domain}
Let $Y_s\sim N(0,2s)$ for fixed $s\ge0$. The actual flow obeys
\begin{equation}\label{eq:heat}
 \E F_t(z+iY_s)=F_{t+s}(z),
 \qquad \E F_t(z+Y_s)=F_{t-s}(z).
\end{equation}
For a complex exponential, $\E e^{qY_s}=e^{sq^2}$. In the half-line
Fourier representation this turns $e^{iu(z+iY_s)}$ into
$e^{iuz+su^2}$. Every interchange is absolutely integrable: integrating
the absolute value first in $Y_s$ is bounded by a constant times
$|K(u)|e^{(t+s)u^2+|\operatorname{Im}z|u}$, integrable by
\eqref{eq:tail}. The opposite-frequency term is handled identically.
For real shifts the absolute-value bound is even smaller, and the
Gaussian characteristic function yields $e^{-su^2}$. This proves both
identities, including their time signs. At $s=0$ the law is a point mass.

\begin{proposition}\label{prop:gaussianlimit}
For every fixed $s\ge0$,
\[
 \E P_n(z+iY_s)\longrightarrow F_s(z)
\]
locally uniformly in $z\in\C$. Each expectation is absolutely integrable.
For each fixed $z,c\in\C$, the same approximants also converge to $F_0$
in the shifted Gaussian $L^1$ norm:
$\E|P_n(z+cY_s)-F_0(z+cY_s)|\to0$.
\end{proposition}
\begin{proof}
Assume $s>0$; $s=0$ is immediate from local uniform convergence.
For a compact disk $|z|\le R$, Lemma~\ref{lem:budget} and
$|z+cy|^2\le2R^2+2|c|^2y^2$ give
\[
 |P_n(z+cy)|\le C_\epsilon e^{2\epsilon R^2}
                                  e^{2\epsilon|c|^2y^2}.
\]
Choose $\epsilon>0$ with $2\epsilon|c|^2<1/(4s)$. This is integrable
under $N(0,2s)$, uniformly in $n$ and in that disk. The limit $F_0$
inherits the same bound. Dominated convergence proves the stated $L^1$
limit at each fixed $z$. To obtain local uniform integral convergence,
first restrict to $|y|\le A$. The affine images of $|z|\le R$ form a
compact set, so the integrand differences tend to zero uniformly there.
Outside $|y|\le A$, the displayed common Gaussian majorant bounds the
integral tail uniformly in $n,z$, and that tail tends to zero as
$A\to\infty$. Finally use \eqref{eq:heat} at $t=0$, $c=i$.
\end{proof}

For polynomials define the finite operator
\begin{equation}\label{eq:operator}
 H_sP=e^{-s\partial_z^2}P
       :=\sum_{j=0}^{\lfloor\deg P/2\rfloor}
                       \frac{(-s)^j}{j!}P^{(2j)}.
\end{equation}
Taylor's finite formula and
$\E Y_s^{2j}=(2j)!s^j/j!$ give exactly
$H_sP(z)=\E P(z+iY_s)$. Thus the locally uniform limit in
Proposition~\ref{prop:gaussianlimit} is the actual finite polynomial
operator limit, with no infinite-dimensional backward-heat series
assumed convergent on an unspecified domain.

The common budget cannot be omitted. For $Y\sim N(0,1)$ the polynomials
$Q_n(z)=1+z^{2n}/(2n-1)!!$ tend to one on every compact complex set,
but $\E Q_n(Y)=2$ for every $n$. The compact convergence follows from
the ratio $R^2/(2n+1)\to0$ of successive nonconstant bounds on $|z|\le R$;
the expectation identity is the exact even Gaussian moment. Local
uniform convergence by itself does not license this Gaussian averaging.
This is a classical diagnostic, not an alternative Xi approximant.

\section{Conjugate-root products contract strips under imaginary shifts}
For a real nonzero polynomial $P$, let
\[
 C_hP(z)=\frac{P(z+ih)+P(z-ih)}2,\qquad h>0.
\]
Its leading coefficient and degree are unchanged, so it is nonzero and
has real coefficients. Suppose all zeros of $P$ obey
$(\operatorname{Im}a)^2\le b^2$.
\begin{lemma}\label{lem:shift}
Every zero of $C_hP$ satisfies
\[
 (\operatorname{Im}z)^2\le\max(b^2-h^2,0).
\]
\end{lemma}
\begin{proof}
If $P$ is constant the assertion is vacuous. Otherwise list its complex
roots with multiplicity in a nonempty multiset $L$, invariant under
conjugation. For $z=x+iy$, direct polynomial expansion gives
\begin{align}\label{eq:pair}
 &|z+ih-a|^2|z+ih-\overline a|^2
       -|z-ih-a|^2|z-ih-\overline a|^2\\
 &\hspace{2em}=8hy\bigl((x-\operatorname{Re}a)^2+y^2+h^2
                                      -(\operatorname{Im}a)^2\bigr).\notag
\end{align}
If $y>0$ and $y^2+h^2>b^2$, this difference is strictly positive for
every root, including real roots and repeated roots. Multiplying the
strict inequalities over $L$, all factors on the larger side are
positive. Conjugation invariance makes each paired product a square
of the unpaired product. Taking nonnegative square roots and restoring
the nonzero leading coefficient gives
\[
 |P(z-ih)|<|P(z+ih)|.
\]
The two terms cannot sum to zero. For a lower-half-plane candidate use
conjugation of the real polynomial $C_hP$. Any nonreal zero violating
the proposed bound has $y^2>0$ and $y^2+h^2>b^2$, so both half-planes
are excluded. Real zeros already satisfy the nonnegative bound.
\end{proof}

Iterating the same shift $n$ times therefore gives
\begin{equation}\label{eq:iterstrip}
 C_h^nP(z)=0\quad\Longrightarrow\quad
       (\operatorname{Im}z)^2\le\max(b^2-nh^2,0).
\end{equation}
The budget identity is
$\max(\max(B,0)-q,0)=\max(B-q,0)$ for $q\ge0$, proved by separating
$B\ge0$ and $B<0$. Thus induction remains valid after the strip has
collapsed. No root is followed through a collision, and no simple-root
or distinct-root assumption is imposed at any stage.

\section{The finite-degree shift limit and closure of zero strips}
Fix $s>0$ and a degree bound $d$. On the finite-dimensional space of
polynomials of degree at most $d$, differentiation $D$ is a bounded
nilpotent operator. Taylor's formula gives the exact finite identity
\[
 C_h=\sum_{j=0}^{\lfloor d/2\rfloor}\frac{(-1)^jh^{2j}}{(2j)!}D^{2j}.
\]
Set $h_n=\sqrt{2s/n}$ for $n\ge1$, and $A=-sD^2$. For any fixed
coefficient-space norm,
\[
 C_{h_n}=I+A/n+E_n,\qquad \|E_n\|\le C/n^2.
\]
There are only finitely many terms; the estimate includes $d<4$, when
the remainder is zero. Also $\|C_{h_n}\|,\|I+A/n\|\le1+C'/n$ after
enlarging the constants. The difference-of-powers telescoping identity
therefore bounds
\[
 \|C_{h_n}^{\,n}-(I+A/n)^n\|
 \le n\frac{C}{n^2}(1+C'/n)^{n-1}\le\frac{Ce^{C'}}n\longrightarrow0.
\]
Finally $A$ is nilpotent, so the binomial expansion of $(I+A/n)^n$
has a fixed finite number of terms. For each fixed $j$,
$\binom nj/n^j\to1/j!$. Hence
\begin{equation}\label{eq:shiftlimit}
 C_{h_n}^{\,n}P\longrightarrow\sum_j\frac{A^jP}{j!}=H_sP
\end{equation}
in coefficients and locally uniformly on the complex plane. This
finite-dimensional proof is a human reorganization of the formal
Hasse-derivative/factorial-jet Euler limit, not a claim that the actual
Xi function belongs to that bounded factorial coefficient domain.

We will use the following precise limit fact twice.
\begin{lemma}\label{lem:closure}
Let nonzero polynomials $Q_n$ have all zeros in
$|\operatorname{Im}z|\le b$, with $b\ge0$. If $Q_n(z)\to f(z)$
pointwise on $\C$, and $f$ is an entire function not identically zero,
then all zeros of $f$ lie in the same strip.
\end{lemma}
\begin{proof}
For a factor $z-a$ with $\operatorname{Im}a\le b$, and real
$b\le y_1\le y_2$, its modulus at $x+iy_1$ is at most its modulus at
$x+iy_2$, since $(y_1-\operatorname{Im}a)^2\le
(y_2-\operatorname{Im}a)^2$. Factorization into all roots with
multiplicity gives
\[
 |Q_n(x+iy_1)|\le|Q_n(x+iy_2)|.
\]
Pass to the pointwise limit. If $f(x+iy_2)=0$ for $y_2>b$, then
$f(x+iy)=0$ for every $y\in[b,y_2]$. The identity theorem makes
$f$ identically zero, a contradiction. Applying the same argument
to $Q_n(-z)$ and $f(-z)$ excludes zeros below $-b$.
\end{proof}
This argument uses the given entireness and nonzero hypothesis of the
limit. Pointwise limits in general need not be entire, and the zero
function cannot be admitted. It avoids assuming convergence of individual
zeros or a continuous ordering of roots.

\begin{proposition}\label{prop:polystrip}
For every nonzero real polynomial $P$, every $s\ge0$, and every bound
$(\operatorname{Im}a)^2\le b^2$ on its initial zeros, the zeros of $H_sP$
obey $(\operatorname{Im}z)^2\le\max(b^2-2s,0)$.
\end{proposition}
\begin{proof}
At $s=0$ the operator is the identity. At $s>0$, the scaled shifts in
\eqref{eq:shiftlimit} all have this same strip by
\eqref{eq:iterstrip}, since $nh_n^2=2s$. Their locally uniform limit
$H_sP$ is nonzero: \eqref{eq:operator} preserves the leading coefficient.
Apply Lemma~\ref{lem:closure} with
$\sqrt{\max(b^2-2s,0)}$ and square the resulting nonnegative bounds.
\end{proof}

\section{Transfer to the actual Xi flow and exact scope}
\begin{proof}[Proof of Theorem~\ref{thm:main}]
The actual polynomials $P_n$ of \eqref{eq:polys} are nonzero real even
polynomials, and \eqref{eq:initialstrip} bounds all their initial zeros
by width one. Proposition~\ref{prop:polystrip} gives all zeros of
$H_tP_n$ the common width $\sqrt{\max(1-2t,0)}$.
Proposition~\ref{prop:gaussianlimit} identifies their locally uniform
limit with the actual $F_t$, after establishing the necessary Gaussian
integrability and common tail bound. That limit is entire and nonzero
by Section 2. Lemma~\ref{lem:closure} transfers the strip, proving the
claimed squared inequality. For $t\ge1/2$ its right side is zero,
forcing $\operatorname{Im}z=0$. The symmetry and positive-origin claims
were already proved directly from the kernel.
\end{proof}

The proof uses two different limiting operations. First, for each fixed
polynomial, increasingly many small symmetric shifts give its finite
heat operator, with degree fixed during that limit. Second, the genuine
divisor-polynomial sequence has increasing degree and converges after
Gaussian averaging because of the shared budget. No degree-uniform
Euler estimate is needed, and no interchange of these two limits is
asserted. This is the useful domain distinction in the formalization.

The coordinate conversion is likewise part of the conclusion. Newman's
classical definition uses $s=1/2+iz$ and a damping parameter in
$e^{-bu^2}$~\cite{Newman}. Here the Xi argument is $z/2$, the initial
strip has width one, and forward $t$ multiplies the arithmetic kernel
by $e^{+tu^2}$. Numerical thresholds cannot be transferred between these
conventions without rescaling both frequency and time. We use only the
bound proved in the displayed normalization.

This manuscript does not prove reality at $t=0$, an exact numerical value of the quadratic
threshold, a sharper arithmetic zero-free region, or any relation between
a higher-even-order threshold and RH. It also proves neither simplicity
of all positive-time zeros nor an unrestricted backward-heat operator on
an entire-function space. Local Hermite splitting and simple-root motion
have their own verified hypotheses and are not needed here. These limits
are mathematical scope statements, not unresolved steps in the proof of
Theorem~\ref{thm:main}.


\section{A finite attained threshold for the actual quadratic flow}
\begin{theorem}\label{thm:finitethreshold}
There is a finite real number $\Lambda\leq1/2$ such that
\[
 \{t\in\R: \text{every zero of }F_t\text{ is real}\}=[\Lambda,\infty).
\]
In particular the endpoint has only real zeros, and for every
$t<\Lambda$ the actual $F_t$ has a nonreal zero.
\end{theorem}
This is the classical finite-threshold conclusion, not a new result.
Newman's Theorem 3~\cite{Newman} proves its antecedent using his
measure classification. We reconstruct instead the frozen library's
actual Gaussian-smoothing and elementary four-point route. In this
normalization $\Lambda$ is exactly the usual de Bruijn--Newman constant:
Rodgers--Tao~\cite{RodgersTao} use precisely our kernel, time and Xi
normalization in their equations (1)--(4). Their separate Theorem 1
proves $\Lambda\geq0$. That deep external result is credited here;
its proof is not reconstructed, formalized anew or used below.

\subsection*{Upper-ray propagation and endpoint closedness}
All the analytic hypotheses of Sections 2--4 hold with any fixed
$F_a$ in place of $F_0$. If every zero of $F_a$ has
$|\operatorname{Im}z|\leq b$, its actual divisor polynomials share
that same strip and common Gaussian budget. Their finite polynomial
heat contraction, actual Gaussian integral limit and vertical strip
closure prove
\[
 F_{a+s}(z)=0\quad\Longrightarrow\quad
 (\operatorname{Im}z)^2\leq\max(b^2-2s,0)\qquad(s\geq0).
\]
These are the same successive limits as before, now at the fixed
starting time $a$; no semigroup on an unrestricted entire-function
space is needed. Taking $b=0$ proves that once all zeros are real,
they remain real at every later time. Thus the set $\mathcal R$
displayed in Theorem~\ref{thm:finitethreshold} is an upper set,
and Theorem~\ref{thm:main} makes it nonempty.

It is also closed. Suppose $a_n\in\mathcal R$ and $a_n\to a$.
For a fixed $n$, every divisor polynomial of $F_{a_n}$ is real rooted.
Its modulus on each line $x+iy$ is nondecreasing for $y\geq0$, by
the elementary real-root factor inequality used in Lemma 6.1.
Passing through its compact limit gives
\[
 |F_{a_n}(x+iy_1)|\leq|F_{a_n}(x+iy_2)|
 \qquad(0\leq y_1\leq y_2).
\]
Time continuity at the two fixed spatial points passes this inequality
to $F_a$. If $F_a(x+iy)=0$ with $y>0$, every point $x+iv$,
$0\leq v\leq y$, would be a zero. The identity theorem contradicts
$F_a(0)>0$. Conjugation excludes a zero below the real axis.
Hence $a\in\mathcal R$. What remains is to exclude
$\mathcal R=\R$ by an actual negative-time obstruction.

\begin{proposition}\label{prop:quadraticnegative}
For some strictly negative $t$, the actual $F_t$ has a nonreal zero.
No numerical time or particular zero is specified by this argument.
\end{proposition}
\subsection*{The elementary four-point inequality}
For $c\geq0$ let $\ell_c(v)=\log(1+cv)$ on $v\geq0$, and define
\[
 D_c=2\log(1+c/2)-\log(1+c).
\]
For every $y\geq1$,
\begin{equation}\label{eq:fourpointfactor}
 \ell_c(y)\geq \ell_c(1)+(y-1)(\ell_c(2)-\ell_c(1))
                         -2(y-1)^2D_c.
\end{equation}
To verify it, put $u=c/(1+c)$ and $v=y-1$. Both are nonnegative.
The scalar inequality $\log(1+w)\geq w-w^2/2$ for $w\geq0$ follows
by differentiating the difference, whose derivative is $w^2/(1+w)$.
Furthermore
\[
 D_c=\log\frac{(1+c/2)^2}{1+c}
 \geq1-\frac{1+c}{(1+c/2)^2}
 =\left(\frac c{c+2}\right)^2,
 \qquad u\leq\frac{2c}{c+2}.
\]
Here $\log r\geq1-r^{-1}$ for $r>0$ follows from
$\log x\leq x-1$ applied to $x=r^{-1}$. Consequently $u^2\leq4D_c$.
Now $\ell_c(y)-\ell_c(1)=\log(1+uv)$ and
$\ell_c(2)-\ell_c(1)=\log(1+u)\leq u$. Therefore
\[
 \log(1+uv)\geq uv-\tfrac12u^2v^2
                 \geq v\log(1+u)-2v^2D_c,
\]
which proves \eqref{eq:fourpointfactor}.

It follows by summation that every profile
$q(v)=A-Bv+\sum_j\alpha_j\log(1+c_jv)$, with $c_j,\alpha_j\geq0$,
satisfies
\begin{equation}\label{eq:fourpointprofile}
 q(y)\geq q(1)+(y-1)(q(2)-q(1))
       -2(y-1)^2\bigl(2q(1/2)-q(0)-q(1)\bigr),\qquad y\geq1.
\end{equation}
Constant and linear functions satisfy this with equality, so adding
$A-Bv$ changes no curvature term. For a nonzero polynomial $P$ with
all roots real and $P(0)\ne0$, its complete real-root factorization gives
\[
 \log|P(iC\sqrt v)|
 =A_P+\frac12\sum_{r\in\operatorname{roots}(P)}
                    \log\bigl(1+(C/r)^2v\bigr).
\]
All roots are counted with multiplicity; each is nonzero. Hence
$A-Bv+\log|P(iC\sqrt v)|$ satisfies \eqref{eq:fourpointprofile}.
If an entire $f$ in our growth class has only real zeros and $f(0)\ne0$,
Section 3 gives such real-rooted approximants with the exact disk
zeros and preserved origin. At every point on the imaginary axis
$f$ is nonzero: a purely imaginary real zero must be zero itself.
Thus continuity of $\log|\cdot|$ passes their inequalities to
$A-Bv+\log|f(iC\sqrt v)|$ at the five fixed values
$0,1/2,1,2,y$. This step needs no differentiation or infinite logarithm sum.

\subsection*{An actual density smoothing identity}
Let $F(x)=K(|x|)$ and $Y\sim N(0,2)$. This is a bounded, strictly
positive continuous even function. For $\epsilon>0$ put
$S_\epsilon(x)=\mathbb E F(x+\epsilon Y)$. These are genuine positive
integrals. Bounded dominated convergence gives
$S_\epsilon(x)\to F(x)$ as $\epsilon\downarrow0$.
The Gaussian variance is exactly $2\epsilon^2$. Pairing the two
half-lines and completing the square gives
\begin{equation}\label{eq:quadraticdensityheat}
 S_\epsilon(x)=\frac{1}{\sqrt\pi\epsilon}
 e^{-x^2/(4\epsilon^2)}
 F_{-1/(4\epsilon^2)}\left(\frac{ix}{2\epsilon^2}\right).
\end{equation}
Indeed its full-line density is
$(2\sqrt\pi\epsilon)^{-1}e^{-(u-x)^2/(4\epsilon^2)}$.
The sum at $u$ and $-u$, on $u\geq0$, equals
\[
 (\sqrt\pi\epsilon)^{-1}e^{-x^2/(4\epsilon^2)}
 e^{-u^2/(4\epsilon^2)}\cosh(ux/(2\epsilon^2)).
\]
Multiplying by $K(u)$ and using $\cos(iw)=\cosh w$ proves the formula.
The bounded-density Gaussian integral is absolutely integrable;
the half-line cosh integral is absolutely integrable by theta decay.
In particular the rightmost transform value is strictly positive real,
so every logarithm used below is legitimate.

\begin{proof}[Proof of Proposition~\ref{prop:quadraticnegative}]
Suppose instead that $F_t$ has only real zeros for every $t<0$.
For each $\epsilon>0$ take $f=F_{-1/(4\epsilon^2)}$ in the preceding
polynomial-limit inequality. Taking logarithms in
\eqref{eq:quadraticdensityheat} shows that
$q_\epsilon(v)=\log S_\epsilon(\sqrt v)$ is exactly one such profile
with
\[
 A=\log(1/(\sqrt\pi\epsilon)),\qquad
 B=1/(4\epsilon^2),\qquad C=1/(2\epsilon^2).
\]
Consequently \eqref{eq:fourpointprofile} holds for $q_\epsilon$.
Fix $y\geq1$ and send $\epsilon\downarrow0$. Positivity of $F$ and
continuity of the real logarithm imply convergence at all five required
points. Thus the same inequality holds for
$q(v)=\log K(\sqrt v)$.

Its right side is one fixed quadratic polynomial $p(y)=ay^2+by+c$.
For $y\geq1$,
$p(y)\geq-(|a|+|b|+|c|)y^2$. Put $L=|a|+|b|+|c|$.
Then $\log K(x)\geq-Lx^4$ for every $x\geq1$, equivalently
\[
 e^{Lx^4}K(x)\geq1\qquad(x\geq1).
\]
But the actual theta estimate $K(x)\leq C e^{9x-\pi e^{4x}/2}$
gives $e^{Lx^4}K(x)\to0$: the exponential $e^{4x}$ dominates the
fixed polynomial $Lx^4+9x$. This contradiction proves that at least
one of the strictly negative quadratic times has a nonreal zero.
\end{proof}

\begin{proof}[Proof of Theorem~\ref{thm:finitethreshold}]
Let $\tau<0$ and a nonreal zero of $F_\tau$ be supplied by the
proposition. Every $r\in\mathcal R$ must satisfy $r>\tau$:
otherwise upper-set propagation would put $\tau$ in $\mathcal R$.
Thus $\mathcal R$ is nonempty and bounded below. Its infimum
$\Lambda$ is a finite real number. Choose $r_n\in\mathcal R$ with
$\Lambda\leq r_n<\Lambda+1/n$ for $n\geq1$, using the defining
infimum property. Closedness gives $\Lambda\in\mathcal R$;
upper propagation then proves $\mathcal R=[\Lambda,\infty)$.
Because $1/2\in\mathcal R$, one has $\Lambda\leq1/2$.
For $t<\Lambda$, nonmembership means exactly that an actual
nonreal zero exists. These statements include the endpoint and
make no assumption about the simplicity of its zeros.
\end{proof}

The named formal endpoint is
\par\noindent\code{RiemannQuadraticFiniteThreshold.real_zeros_iff},
with the finite lower-bound obstruction in that same module and the
smoothing identity in \code{RiemannQuadraticDensityHeat}.
The proof here first produces a strictly negative nonreal time; the
formal route first excludes all-realness at every real time and then
uses its attained threshold to produce such a negative time.
This explicitly labelled human reordering proves the same conclusions.

By the exact Xi normalization, $\Lambda\leq0$ is the standard
threshold formulation of RH. This paper proves neither that bound nor
the numerical value of $\Lambda$. The separate published lower bound
$\Lambda\geq0$ is not an open sign question in the literature and
is not a consequence of the elementary proof above. The bound $1/2$
is retained solely as the classical bound reconstructed here, without
claiming it is the best available bound. No quartic or sixth-order
threshold is identified with $\Lambda$.

\section{Formal correspondence, primary comparison and verification}
The map supplied with this paper records exact declarations, source
hashes and line links at the fixed baseline. In the source,
$F_t$ is \code{RiemannEvenFlow.flow 1 t} and $H_t$ is
\code{PureOrderPolynomialTransport.transport 2 (-t)}.
The main chain is:
\begin{itemize}
\item \code{RiemannQuadraticZeroProductTails.zero_time_strip} derives
width one from the actual zeta nonvanishing theorem for
$\operatorname{Re}s>1$, without an external $7/8$ premise.
\item \code{RiemannQuadraticPolynomialHeat} supplies the actual even
polynomials, disk-zero retention, normalization and common Gaussian budget.
\item \code{RiemannQuadraticUniformHeat} and
\code{RiemannQuadraticOperatorApproximation} identify genuine locally
uniform Gaussian and polynomial-operator limits.
\item \code{SymmetricPolynomialStrip},
\code{SymmetricPolynomialTaylor} and
\code{SymmetricPolynomialHeatLimit} retain multiplicities, the scaled
shift budget and the fixed-degree heat limit.
\item \code{PolynomialStripLimit} proves strip closure and finite
polynomial heat contraction.
\item \code{RiemannQuadraticGlobalStrip.flow_zero_strip} and
\code{real_zeros_after_half} give exactly Theorem~\ref{thm:main}.
\end{itemize}
The proof above uses classical Hadamard factorization with Jensen's
estimate for the divisor-polynomial construction, and a finite-dimensional
coefficient telescoping proof for the Euler step. These are explicit human
presentations of the named formal conclusions, rather than shorter new
Lean declarations. The source's annular divisor-tail and factorial-jet
proofs remain separately mapped. Fundamental complex-analysis inputs
(identity theorem, Jensen and Hadamard), the elementary theta identity,
Mellin continuation and the Euler-product zero-free half-plane are stated
where used. No new Lean proof or research-source modification is made.

The classical status of the main bound is settled by the exact primary
antecedent: Newman's 1976 Theorem 1 states the broader strip contraction
and credits de Bruijn. Its complete statement was inspected, including
its decay and evenness hypotheses and its growing quadratic multiplier.
Its strict initial strip can be widened from width one to $1+\epsilon$
and then $\epsilon\downarrow0$, giving the closed-boundary convention
used here.
The original de Bruijn full-text comparison remains incomplete in this
editorial audit; the source's citation identifies its finite-shift route,
but that is not presented as a fresh full-text review by the editor.
McMullen's primary course Hadamard passages provide the foundational
factorization context. No novelty is claimed for any general strip,
Gaussian heat or entire-product mechanism. The contribution of this
companion is a readable, exact-domain reconstruction of an actual verified
arithmetic proof chain. A reviewed draft is not an external submission.

The frozen baseline already has a complete successful Lean root build
and transitive axiom audit, with the standard classical axioms only.
Publication work checks those recorded evidence hashes and all mapped
source bytes; it does not consume a redundant research build. Independent
finite algebra checks, when supplied, verify only the displayed pair
identity and finite heat/shift transcriptions, not the analytic theorems.
The source/PDF compilation and visual-review receipts are kept separately.

\begin{thebibliography}{9}
\bibitem{deBruijn} N. G. de Bruijn,
\emph{The roots of trigonometric integrals}, Duke Math. J. \textbf{17}
(1950), 197--226. \url{https://doi.org/10.1215/S0012-7094-50-01720-0}.
Original full-text comparison remains incomplete.
\bibitem{Newman} C. M. Newman,
\emph{Fourier transforms with only real zeros}, Proc. Amer. Math. Soc.
\textbf{61} (1976), no.~2, 245--251. Theorem 1, p.~246.
\url{https://sites.math.northwestern.edu/~auffing/papers/Newman.pdf}.
\bibitem{McMullen} C. McMullen,
\emph{Complex Analysis}, Harvard Math 213a course notes, 2023,
Theorems 3.15 and 3.18, printed pp.~80--82.
\url{https://people.math.harvard.edu/~ctm/home/text/class/harvard/213a/23/html/home/course/course.pdf}.
\bibitem{RodgersTao} B. Rodgers and T. Tao,
\emph{The de Bruijn--Newman constant is non-negative},
author manuscript v5, July 3, 2021, equations (1)--(4) and Theorem 1,
printed pp.~2--3. \url{https://arxiv.org/abs/1801.05914}.
The external nonnegativity proof is not reproduced here.
\end{thebibliography}
\end{document}
